Flood susceptibility maps are judged by how well they rank the ground that floods, yet they are rarely compared with the simplest alternative, the record of where floods have already been. This study makes that comparison for infrastructure mapped in OpenStreetMap in three regions of Bangladesh, using flood extents mapped from Sentinel-1 radar for thirteen events as labels and reference, and scoring every model on 10\,km blocks held out from training over \RRBenchSeeds{} block assignments. Trained on floods up to 2022, a gradient-boosted model ranked the ground the 2024 floods reached at an AUC of \RRTempAUC{} in Rangpur and Rajshahi and \SylhetTempAUC{} in Sylhet, below the record of earlier flooding at \RRTempPastAUC{} and \SylhetTempPastAUC{}. Given that record as a feature, the model beat both, at \RRPastFeatAUC{} and \SylhetPastFeatAUC{}, because terrain still ranks the ground no earlier flood reached. On the south-west coast, where most assets flooded by the 2024 cyclone stood on ground that had not flooded before, the model beat the record, at \CoastalTempAUC{} against \CoastalTempPastAUC{}. Training on radar-observed extents rather than terrain-threshold labels raised the AUC by \RRLabelGap{}, \SylhetLabelGap{} and \CoastalLabelGap{} on held-out areas and by \RRTempProxyGap{}, \SylhetTempProxyGap{} and \CoastalTempProxyGap{} on the 2024 floods. A graph neural network linking neighbouring assets added nothing over tabular models, and elevation carried almost all of the terrain signal. The radar masks agreed only moderately with an independent MODIS flood map, so the scores measure agreement with flooding as radar sees it. Susceptibility studies with a radar flood history should report that record as a baseline.
